\documentclass[10pt,a4paper]{amsart}
\usepackage[margin=19mm]{geometry}
\usepackage{amsmath,amssymb,mathtools}
\usepackage[scr=rsfs,cal=euler]{mathalfa}
\usepackage{xfrac}
\usepackage[unicode,hidelinks,bookmarks=false]{hyperref}
\newcommand{\ie}{i.e.\ }
\newcommand{\cf}{cf.\ }
\newcommand{\resp}{respectively\ }
\newcommand{\Subsection}{\subsection}
\providecommand{\nobreakdash}{\nobreak\hbox{-}}
\setlength{\emergencystretch}{2em}
\multlinegap0pt
\usepackage{mathtools}
\usepackage{amssymb}
\usepackage{multirow}
\usepackage{tikz}
\usepackage[shortlabels]{enumitem}
\usepackage{booktabs}
\usepackage{float}
\usepackage{dsfont}
\usepackage{bm}
\usepackage{bbm}
\newtheorem{theorem}{Theorem}
\newtheorem{proposition}{Proposition}
\newtheorem{claim}{Claim}
\newtheorem{poc}{Poc}
\title{\bf\Large Rational codegree Tur\'an density of hypergraphs}
\author{Jun Gao, Oleg Pikhurko, Mingyuan Rong, Shumin Sun}
\date{Source: arXiv:2601.00758v1 (CC BY 4.0)}
\begin{document}
\maketitle

\noindent\emph{Source and licence.} \bf\Large Rational codegree Tur\'an density of hypergraphs. Version arXiv:2601.00758v1 (CC BY 4.0), distributed by arXiv under the Creative Commons Attribution 4.0 International licence, http://creativecommons.org/licenses/by/4.0/, as recorded in the arXiv licence metadata for this exact version and shown on the abstract page https://arxiv.org/abs/2601.00758v1. Full licence notice: \texttt{licenses/arxiv-2601.00758-CC-BY-4.0.txt}.\par\smallskip

\noindent\textit{Editorial scope.} Counted unit: the article's theorem thm:F\_r (source0.tex lines 486--488). The article's complete original proof of it is delivered with it (source0.tex lines 489--555, one \texttt{proof} environment of 67 lines). No mathematics is written, repaired or re-derived: the statement and the proof are the article's own lines, in the article's own notation, and no retained line is substituted or re-flowed.  Retained uncounted as the source-local context the proof needs: lines 216--220; lines 453--454.  Nothing was omitted from any retained span, so the package carries no omission record.  No retained line contains a citation, so the wrapper carries no bibliography and no bibliography entry.  No retained line contains a figure environment, an image-inclusion command or drawing code, so no image is delivered and the package declares no asset dependency.  Editorial formatting only, in addition: an AMS \texttt{amsart} preamble that restates the article's own declarations and macros used by the retained lines, namely the environments \texttt{theorem}, \texttt{proposition}, \texttt{claim}, \texttt{poc} on separate counters, together with the standard packages those lines need.  The retained environments print in this document as Theorem 1, Proposition 1, Claim 1, Poc 1 in order of appearance, whereas the article prints the same items with the numbering of its own full document.  The complete native source of the article as distributed in the arXiv e-print for this exact version is bundled unchanged as \texttt{sources/192024/source0.tex}.
 \begin{theorem}\label{thm:non-principal}
  For every $k\ge 3$, there exist $k$-graphs $F_1$ and $F_2$ such that
     \[0<\gamma(\{F_1, F_2\})<\min\{\gamma(F_1),\gamma(F_2)\}.
     \]
 \end{theorem}

\begin{proposition}\label{ob:KZ} If a $k$-graph $F$ is not $r$-colourable, then $\gamma(F)\ge \frac{r-1}{r}$.
\end{proposition}

\begin{theorem}\label{thm:F_r}
    Fix any $k \ge 3$. Then, for every $r\ge1$, the $k$-graph $F^k_r$ is not $r$-colourable and satisfies $\gamma(F^k_r) = \frac{r-1}{r}$.
\end{theorem}

\begin{proof}
We prove the theorem by induction on~$r$.
Although the case $r=1$ is trivial and the induction step handles the case $r=2$ correctly, we present a separate proof for $r=2$ (since $F_2^k$ plays a crucial role in Theorem~\ref{thm:non-principal}).
%, we first prove that $F^k_2$ is not 2-colourable and $\gamma(F^k_2)=1/2$. 
\begin{claim}\label{cl:gammaF2}
    The $k$-graph $F^k_2$ is not 2-colourable and $\gamma(F^k_2)=1/2.$
\end{claim}
\begin{poc}
    Let us first show that $F^k_2$ is not $2$-colourable. Suppose, for contradiction, that $F^k_2$ is $2$-colourable, and let $V(F^k_2)=A\cup B$ be the partition such that neither $A$ nor $B$ contains an edge. Recall that $V(F^k_2)= S\cup \bigcup_{i\in [m]}V(\tilde{X}_{i})$, where $|S|=2k-3$. By the pigeonhole principle, there exists a set $S'\subseteq S$ of size $k-1$ entirely lying in $A$ or $B$. Without loss of generality, assume that $S'\subseteq A$. Let $i\in [m]$ be such that $S'=X_i$. Consider the corresponding set $\tilde{X}_{i}$. Since $\tilde{X}_{i}$ is actually an edge of $F^k_2$, it must have at least one vertex $v\in \tilde{X}_{i}$ lying in $A$; otherwise, $\tilde{X}_{i}\subseteq B$ contradicting the assumption that $B$ is edge-free. But then $X_i\cup \{v\}$ is an edge lying in $A$, a contradiction. Hence,  $F^k_2$ is not $2$-colourable.

By Proposition~\ref{ob:KZ}, we have that $\gamma(F^k_2)\ge 1/2$. Let us show the upper bound $\gamma(F^k_2)\le 1/2$. Let $\varepsilon>0$ be an arbitrarily small constant, and assume that $n$ is sufficiently large. We claim that any $k$-graph $H$ with minimum codegree $\delta_{k-1}(H)\ge (1/2+\varepsilon)n$ must contain $F^k_2$ as a subgraph.
        
Identify  an arbitrary $(2k-3)$-subset of $V(H)$ with $S$. Its $(k-1)$-subsets are labelled by $X_1,X_2,\dots,X_m$, where $m\coloneqq \binom{2k-3}{k-1}$. We find suitable sets  $\tilde{X}_1,\dots,\tilde{X}_m\subseteq V(H)$ one by one. Suppose that $i\ge 1$ and we have already defined $\tilde{X}_1,\dots,\tilde{X}_{i-1}$. Let 
\[
A:=N(X_i)\setminus\Big(S\cup \bigcup_{j=1}^{k-1} \tilde{X}_j\Big).
\]By the codegree condition of $H$, we have that
    $$
    |A|\ge d(X_i)-|S|-(i-1)k\ge \left(\frac{1}{2}+\varepsilon\right)n-(2k-3)-(i-1)k\ge %\left(\frac{1}{2}+\frac{\varepsilon}{2}\right)n\ge
     k-1.
    $$
    Take an arbitrary $(k-1)$-subset $Y$ of $A$. Since 
    \[
    d(Y)\ge \left(\frac12+\varepsilon\right)n> |V(H)|-|A|,
    \] 
    there must exist a vertex $y\in N(Y)\cap A$. We take $Y\cup \{y\}\in E(H)$ for $\tilde{X}_i$.  After $m$ steps, the sets $S, \tilde{X}_{1},\dots,\tilde{X}_{m}$ form a copy of $F^k_2$ in $H$, finishing the proof of the claim.  
\end{poc}

In the induction step for $r\ge 3$, we have to show that  $F^k_{r}$ is not $r$-colourable and $\gamma(F^k_{r}) =\frac{r-1}{r}$. The proof is very similar to that of Claim~\ref{cl:gammaF2}. 
%For completeness, we include the proof.
First, we show that $F^k_r$ is not $r$-colourable.  Suppose, for contradiction, that $F^k_r$ is $r$-colourable. Then $V(F^k_r)$ can be partitioned into $r$ parts $A_1,A_2,\dots,A_r$, each of which induces no edge. Recall that \[V(F^k_r) \coloneqq S\cup \bigcup^m_{j=1}V\left(\tilde{X}_j\right) \text{ and } |S| = (k-2)r+1.\]
By the pigeonhole principle, there exists a set $S' \subseteq S$ of size $k-1$ entirely lying in some part~$A_i$.  Without loss of generality, assume that $S'\subseteq A_r$. Let $S'=X_j$. Consider the corresponding copy of $F^k_{r-1}$, namely $\tilde{X}_j$.
As every vertex of $\tilde{X}_j$ forms an edge together with $X_j$, we must have that
\[
\tilde{X}_j \cap A_r = \varnothing,
\]
which contradicts the fact that $\tilde{X}_j$ is not $(r-1)$-colourable.
Therefore, $F^k_r$ is not $r$-colourable.


Next, we show that $\gamma(F^k_r)=\frac{r-1}{r}$. By Observation~\ref{ob:KZ}, it remains to prove that $\gamma(F^k_r)\le \frac{r-1}{r}$. 
Given any (arbitrarily small) constant $\varepsilon>0$, let $n$ be sufficiently large. Let us show that any $k$-graph $H$ with minimum codegree $\delta_{k-1}(H)\ge (\frac{r-1}{r}+\varepsilon)n$ must contain $F^k_r$ as a subgraph.

Take an arbitrary vertex set $S \subseteq V(H)$ of size $r(k-2)+1$, and label its $(k-1)$-subsets by $X_1,X_2,\dots,X_m$. Our goal is to iteratively find, for each $X_i$, a copy of $F^k_{r-1}$ in $N(X_{i})$, denoted by $\tilde{X_i}$, with the additional requirement that $\tilde{X_1},\dots,\tilde{X_m}$ and $S$ are pairwise disjoint. Suppose that we have obtained $\tilde{X}_1,\dots,\tilde{X}_{i-1}$ for some $i\in[m]$.
We now have to find $\tilde{X}_{i}$.
Define
\[
H_i \coloneqq  H\Big[\, N(X_{i}) \setminus \Big(S \cup \bigcup_{j\in[i-1]} \tilde{X}_j\Big) \,\Big].
\]
Then we have
\[
|V(H_i)|
\ge \left(\frac{r-1}{r}+\varepsilon\right)n  - (r(k-2)+1)
- (i-1)\,|V(F_{r-1}^k)|
\ge \left(\frac{r-1}{r}+\frac{\varepsilon}{2}\right)n,
\]
where the last inequality holds since $n$ is  sufficiently large.
Consequently,
\[
\delta_{k-1}(H_i)
\ge \delta_{k-1}(H) - \bigl(n-|V(H_i)|\bigr)
\ge \left(\frac{r-2}{r-1}+\varepsilon'\right)|V(H_i)|,
\]
for some constant $\varepsilon'>0$ (depending on $\varepsilon$ and $r$ only).
By the induction hypothesis, $\gamma(F^k_{r-1})=\frac{r-2}{r-1}$.
Hence, $H_i$ contains a copy of $F^k_{r-1}$, which we denote by $\tilde{X}_i$.
After we have processed all $i\in[m]$, the sets $S, \tilde{X}_1,\dots,\tilde{X}_m$ form a copy of $F^k_r$ in $H$. This finishes the proof of the theorem.
\end{proof}

\end{document}
